% -----------------------------------------------------------------------------
% Tipografia e idioma
% -----------------------------------------------------------------------------
\usepackage{fontspec}
\setmainfont{Liberation Serif}
\setsansfont{Liberation Sans}
\setmonofont{Liberation Mono}
\usepackage[brazil]{babel}
\usepackage{microtype}

% -----------------------------------------------------------------------------
% Layout principal: UFRJ/SiBI 2025, documento digital em página única.
% A4; margens esquerda/superior 3 cm e direita/inferior 2 cm.
% -----------------------------------------------------------------------------
\usepackage[
  a4paper,
  left=3cm,
  top=3cm,
  right=2cm,
  bottom=2cm,
  headheight=15pt,
  headsep=0.55cm,
  includehead=false,
  includefoot=false
]{geometry}

\usepackage{setspace}
\OnehalfSpacing
\setlength{\parindent}{1.25cm}
\setlength{\parskip}{0pt}
\setlength{\emergencystretch}{2em}
\clubpenalty=10000
\widowpenalty=10000
\displaywidowpenalty=10000

% -----------------------------------------------------------------------------
% Cabeçalho: apenas número da página no canto superior direito.
% Usa a interface nativa da classe memoir/abnTeX2.
% -----------------------------------------------------------------------------
\makepagestyle{coinfosim}
\makeoddhead{coinfosim}{}{}{\small\thepage}
\makeevenhead{coinfosim}{\small\thepage}{}{}
\makeoddfoot{coinfosim}{}{}{}
\makeevenfoot{coinfosim}{}{}{}
\makeheadrule{coinfosim}{\textwidth}{0pt}
\aliaspagestyle{plain}{coinfosim}

% -----------------------------------------------------------------------------
% Títulos: hierarquia UFRJ com influência visual do template SBC.
% -----------------------------------------------------------------------------
\renewcommand{\ABNTEXchapterfont}{\bfseries\sffamily}
\renewcommand{\ABNTEXchapterfontsize}{\fontsize{13}{15}\selectfont}
\renewcommand{\ABNTEXsectionfont}{\bfseries\sffamily}
\renewcommand{\ABNTEXsectionfontsize}{\fontsize{12}{14}\selectfont}
\renewcommand{\ABNTEXsubsectionfont}{\bfseries\sffamily}
\renewcommand{\ABNTEXsubsectionfontsize}{\fontsize{12}{14}\selectfont}
\renewcommand{\ABNTEXsubsubsectionfont}{\itshape\sffamily}
\renewcommand{\ABNTEXsubsubsectionfontsize}{\fontsize{12}{14}\selectfont}
\setlength{\beforechapskip}{-1.2cm}
\setlength{\afterchapskip}{1.5\baselineskip}
\setpnumwidth{2em}
\setrmarg{3em}

% -----------------------------------------------------------------------------
% Matemática, tabelas e gráficos
% -----------------------------------------------------------------------------
\usepackage{amsmath,amssymb,mathtools,bm}
\usepackage{booktabs}
\usepackage{tabularx}
\usepackage{longtable}
\usepackage{array}
\usepackage{multirow}
\usepackage{siunitx}
\sisetup{output-decimal-marker={,}}
\usepackage{graphicx}
\usepackage{adjustbox}
\usepackage{pdflscape}
\usepackage{subcaption}
\usepackage{float}

% -----------------------------------------------------------------------------
% Legendas: identificação acima; fonte abaixo. Fonte sans 10 pt conforme SBC.
% -----------------------------------------------------------------------------
\usepackage{caption}
\captionsetup{
  position=above,
  font={sf,small},
  labelfont=bf,
  textfont=bf,
  margin=0.8cm,
  justification=justified,
  singlelinecheck=true,
  skip=6pt
}
\captionsetup[subfigure]{
  position=below,
  font={sf,footnotesize},
  labelfont=bf,
  textfont=normalfont,
  margin=0pt,
  justification=centering
}

% -----------------------------------------------------------------------------
% Referências e URLs: sistema autor-data ABNT, processado por Biber.
% -----------------------------------------------------------------------------
\usepackage{csquotes}
\usepackage[
  backend=biber,
  style=abnt,
  ittitles,
  giveninits=true,
  repeatfields=true,
  justify
]{biblatex}
\addbibresource{references/references.bib}
\AtBeginBibliography{\SingleSpacing\raggedright}
\usepackage{xurl}
\hypersetup{
  pdftitle={\ReportTitle: \ReportSubtitle},
  pdfauthor={\ReportAuthor},
  pdfsubject={Relatório acadêmico do projeto CoInfoSim},
  pdfkeywords={canais de informação, cooperação preditiva, perfil de cooperação preditiva, dados sintéticos, seleção de subconjuntos, simulação de Monte Carlo},
  bookmarksopen=true,
  bookmarksnumbered=true
}

% -----------------------------------------------------------------------------
% Identidade visual: uma cor de destaque sóbria e neutros em cinza.
% A paleta preserva legibilidade em impressão em tons de cinza.
% -----------------------------------------------------------------------------
\usepackage[most]{tcolorbox}
\usepackage{xcolor}
\definecolor{CoinAccent}{RGB}{25,80,140}       % azul institucional (destaque único)
\definecolor{CoinAccentDark}{RGB}{18,55,95}    % variação para títulos
\definecolor{CoinAccentPale}{RGB}{237,242,248} % fundo claro de caixas de definição
\definecolor{CoinGray}{RGB}{95,95,95}          % cinza médio para molduras discretas
\definecolor{CoinGrayPale}{RGB}{246,246,246}   % fundo claro neutro
\definecolor{DraftGray}{RGB}{245,245,245}
\definecolor{DraftBorder}{RGB}{120,120,120}
\definecolor{LinkBlue}{RGB}{25,80,140}
\hypersetup{colorlinks=true,linkcolor=black,citecolor=black,urlcolor=LinkBlue}

% Destaque sutil nos títulos de capítulo e seção (imprime quase preto em cinza).
\renewcommand{\ABNTEXchapterfont}{\bfseries\sffamily\color{CoinAccentDark}}
\renewcommand{\ABNTEXsectionfont}{\bfseries\sffamily\color{CoinAccentDark}}
\renewcommand{\ABNTEXsubsectionfont}{\bfseries\sffamily}
\renewcommand{\ABNTEXsubsubsectionfont}{\itshape\sffamily}

% Rótulos de legenda na cor de destaque; texto permanece neutro.
\captionsetup{labelfont={bf,color=CoinAccentDark}}
\captionsetup[subfigure]{labelfont={bf,color=CoinAccentDark}}

\newtcolorbox{editorialbox}[1][]{
  breakable,
  colback=DraftGray,
  colframe=DraftBorder,
  boxrule=0.4pt,
  arc=1mm,
  left=3mm,right=3mm,top=2mm,bottom=2mm,
  fonttitle=\bfseries\sffamily,
  title={Nota editorial},
  #1
}

% Listas compactas apenas quando necessário.
\usepackage{enumitem}
\setlist[itemize]{topsep=3pt,itemsep=2pt,parsep=0pt}
\setlist[enumerate]{topsep=3pt,itemsep=2pt,parsep=0pt}
